\chapter{Hasard et probabilités}
\label{corriges:11}
\section*{Exercice 11.1 -- Une urne}
Il y a $10$ boules équiprobables. $P(R)=3/10$, $P(B)=2/10=1/5$, $P(V)=5/10=1/2$. Le contraire de vert a pour probabilité $1-1/2=1/2$. Les couleurs ne sont pas équiprobables, puisqu'elles correspondent à des nombres de boules différents.

\section*{Exercice 11.2 -- Deux pièces}
Les quatre issues équiprobables sont $PP,PF,FP,FF$. Exactement un pile : $PF,FP$, soit $1/2$. Au moins un pile : $PP,PF,FP$, soit $3/4$ ; c'est aussi $1-P(FF)$. Deux fois le même côté : $PP,FF$, soit $1/2$. Regrouper « un pile et une face » en une seule issue sans tenir compte de l'ordre ferait perdre son double poids.

\section*{Exercice 11.3 -- Avec puis sans remise}
Avec remise, les neuf couples de boules sont équiprobables ; quatre sont rouges :
\[(R_1,R_1),\quad(R_1,R_2),\quad(R_2,R_1),\quad(R_2,R_2).\]
Donc $P(RR)=4/9$.
Sans remise, les couples de même boule sont exclus : six couples restent, dont deux rouges. Donc $P(RR)=2/6=1/3$. Sur un arbre, les produits correspondants sont $(2/3)(2/3)$ puis $(2/3)(1/2)$.

\section*{Exercice 11.4 -- Deux dés et simulation}
Les couples favorables sont $(3;6),(4;5),(5;4),(6;3)$, donc $P=4/36=1/9\approx0{,}1111$. La fréquence observée vaut $114/1000=0{,}114$, proche de cette valeur mais non exactement égale. Les fluctuations sont normales. Un entier uniforme de $2$ à $12$ rendrait les onze sommes équiprobables, alors que leurs nombres de couples diffèrent.
